% !TEX root = ../main.tex
\setcounter{chapter}{-1}
\chapter{仕様・変更・証明の地図}
\label{chap:ch00_map}
\BookIndex{ただしさ}{正しさ}
\BookIndex{しよう}{仕様}
\BookIndex{けんろん}{圏論}
\BookIndex{Lean 4}{Lean 4}

\begin{chapterabstract}
本章は、本書全体の地図です。
圏論や Lean 4 の細部に入る前に、ソフトウェア開発で何を「正しい」と呼ぶのか、変更前後で何を保存したいのか、そして正しさと保存についての主張をどのように検査可能な形へ落とすのかを整理します。
テスト、型、仕様、証明は、互いに競合する道具ではありません。
それぞれが異なる粒度で正しさを支えます。
本章では、圏論を合成と変換の言語として、Lean 4 を、仕様を検査する小さな実験室として位置づけます。
\end{chapterabstract}

\begin{learninggoals}
\item ソフトウェアにおける正しさを、実行結果、型、仕様、不変条件、変更前後の意味保存に分けて説明できます。
\item テスト、型、仕様、証明がそれぞれ何を得意とし、何を保証しないのかを区別できます。
\item 圏論の語彙を、型、関数、変換、パイプライン、マイグレーション、adapter と結びつけて読めます。
\item Lean 4 を、本番コードの代替ではなく、仕様の核を検査する小さなモデルとして位置づけられます。
\item 本書で扱うケーススタディが、どのような「保存したい性質」を扱うのかを把握できます。
\end{learninggoals}

\section{本書が扱う問い}

本書は、圏論の定義を順番に暗記する本ではありません。
また、Lean 4 の構文を網羅的に説明する本でもありません。
主題は、ソフトウェアの仕様と変更を、より正確に扱うことです。

日々の開発では、次のような判断が繰り返し現れます。

\begin{itemize}
\item 関数を分割した後も、外から見た結果は変わっていませんか。
\item API レスポンスの形を変えた後も、旧クライアントの期待を壊していませんか。
\item DB スキーマを移行した後も、必要なクエリ結果は保存されていますか。
\item 文字列や JSON を encode して decode したとき、元の情報に戻りますか。
\item 口座残高や注文合計のような不変条件は、操作後にも保たれていますか。
\end{itemize}

これらは別々の作業に見えます。
しかし、どの作業も「何かを別の形へ変換し、その変換が期待した性質を保存するか」を問うています。
この共通部分を取り出すと、圏論と Lean 4 が同じ方向を向いていることが分かります。

圏論は、変換と合成を整理するための言語です。
Lean 4 は、その言語で書いた小さな主張を機械的に確認するための道具です。
本書では、圏論を抽象数学の中だけで扱わず、Lean も学術的な証明だけに限定しません。
どちらも、変更に強い設計を考えるための実務上の補助線として使います。

図\ref{fig:ch00-book-roadmap}は、仕様と変更の問いから出発し、Lean 4、圏論、証明技法を経てケーススタディへ進む本書の順序を示します。

\FigurePlaceholder{fig-ch00-book-roadmap.png}{0.92}{本書のロードマップ}{fig:ch00-book-roadmap}

本章では、まだ本格的な Lean 証明には入りません。
最初に必要なのは、証明を書くことではなく、「何を証明すべきか」を見分けることです。
正しさの種類を混同したまま Lean を使うと、証明は書けても、実務上あまり重要でない性質だけを確認してしまいます。
一方、証明すべき核を適切に切り出せれば、短い定理でも設計レビューに大きな価値を持ちます。

\section{ソフトウェアにおける「正しさ」とは何か}
\label{sec:ch00-correctness}

\subsection{正しさは一種類ではない}

ソフトウェアの正しさは、単に「バグがない」という一語では説明しきれません。
実務では、いくつもの異なる正しさが重なっています。

\begin{table}[htbp]
\centering
\caption{正しさの分類}
\label{tab:ch00-correctness-kinds}
\begin{tabular}{p{0.20\linewidth}p{0.34\linewidth}p{0.36\linewidth}}
\toprule
種類 & 典型的な問い & 例 \\
\midrule
実行結果の正しさ & この入力で期待した出力になるか & 料金計算で、入力 1000 円に対して 1100 円が返る \\
型の正しさ & 不正な形の値を受け取らないか & ユーザー ID と金額を取り違えない \\
仕様の正しさ & どの入力でも期待する性質を満たすか & 出金後も残高が負にならない \\
変更の正しさ & 変更前後で意味が保存されるか & リファクタリング後も計算結果が同じ \\
移行の正しさ & 古い形式と新しい形式の間で情報が保存されるか & \texttt{rollback (migrate u) = u} が成り立つ \\
合成の正しさ & 部品をつないでも約束が保たれるか & adapter を通す順序を変えても業務ロジックの結果が同じ \\
\bottomrule
\end{tabular}
\end{table}

これらの正しさは、互いに置き換えられません。
ある入力に対するテストが通っても、すべての入力で不変条件が保たれるとは限りません。
型が付いていても、金額計算の式が仕様どおりであるとは限りません。
仕様書に「互換性がある」と書いてあっても、互換性の方向と情報損失の有無は別途明確にする必要があります。

「テストがあるから仕様は不要」「型があるから証明は不要」「証明したから実装全体が安全」という言い方は、いずれも保証範囲を超えています。
テスト、型、仕様、証明は、それぞれ異なる不足を補います。
どれか一つがほかを完全に置き換えるわけではありません。

\subsection{正しさは保存性として現れる}

本書では、正しさをしばしば「保存される性質」として捉えます。
この捉え方は、変更を扱うときに特に有効です。

たとえば、関数をリファクタリングするときに保存したいものは、内部構造ではなく外から見た結果です。
DB スキーマを変えるときに保存したいものは、テーブルの形そのものではなく、業務上必要なクエリの意味です。
JSON encoder と decoder で保存したいものは、文字列の見た目ではなく、データが持っていた情報です。

実務で「安全な変更」と呼ぶものの多くは、「すべてを変えない変更」ではありません。
安全な変更とは、変えてよい部分と保存すべき部分を区別し、保存すべき部分が実際に保たれる変更です。

この見方を採ると、レビューで確認すべき問いが変わります。

\begin{itemize}
\item 何が変わりましたか。
\item 何を変えてはいけませんか。
\item 変えてはいけないものは、テスト、型、仕様、証明のどれで確認するのが適切ですか。
\item その確認は、具体例に対するものですか、それともすべての場合に対するものですか。
\end{itemize}

これらの問いが、本書全体の出発点になります。

\section{テスト、型、仕様、証明の役割分担}
\label{sec:ch00-test-type-spec-proof}

\subsection{四つの道具を重ねて使う}

テスト、型、仕様、証明は、同じ「正しさ」を別の角度から支えます。

図\ref{fig:ch00-correctness-layers}は、具体例を確かめるテストから、モデル上の全入力を扱う証明まで、四つの道具が受け持つ範囲を分けて示します。

\FigurePlaceholder{fig-ch00-correctness-layers.png}{0.90}{正しさを支える四つの道具}{fig:ch00-correctness-layers}

四つの道具を混同しないため、本書ではそれぞれを次の意味で使います。
\begin{itemize}
\item \textbf{テスト}：具体的な入力例について期待結果を確認します。
\item \textbf{型}：値や関数がどの形を持つかを表し、不正な組み合わせを事前に排除します。
\item \textbf{仕様}：システムや関数が満たすべき性質を明文化します。
\item \textbf{証明}：その仕様が定義された範囲のすべての場合で成り立つことを示します。
\end{itemize}

この区別を起点に、確認したい性質ごとに適切な道具を選びます。

テストは、外部システムとの接続、実 DB の挙動、UI、ブラウザ差異、パフォーマンス、運用上の例外を確認するために重要です。
これらは、Lean の小さなモデルだけでは扱いにくい対象です。
一方、テストで選べる入力例は有限個です。
入力空間が広いとき、「選んだ例では正しい」ことと「すべての入力で正しい」ことの間には大きな差があります。

型は、テストよりも広い範囲で不正な形を排除します。
たとえば、\texttt{Nat} を使えば負の数はそもそも表現できません。
\texttt{Option User} を使えば、値が存在しない場合を型に表せます。
しかし、型が付いているからといって、関数が業務仕様を満たすとは限りません。
\texttt{Nat \to Nat} という型を持つ関数は無数にあり、その中には料金計算として正しいものも誤ったものもあります。

仕様は、期待する性質を言葉にします。
自然言語の仕様書は人間に読みやすい一方で、曖昧さも持ちます。
「適切に処理する」「互換性を保つ」「一貫性を維持する」という表現だけでは、機械は検査できません。
そのため、仕様の一部を型、関数、述語、定理へ分解する必要があります。

証明は、仕様として切り出した性質が定義された範囲で常に成り立つことを示します。
ただし、証明が保証するのはモデルの内側だけです。
Lean で \texttt{UserV1} と \texttt{UserV2} という小さなモデルの round-trip を証明しても、本番のシリアライザ、DB 制約、文字コード、外部 API の全挙動まで自動的に保証されるわけではありません。

\subsection{どの道具を使うべきか}

実務では、すべてを証明しようとするのではなく、性質に合う道具を選ぶ必要があります。
次の表は、その判断の出発点です。

\begin{table}[htbp]
\centering
\caption{性質に応じた道具の選択}
\label{tab:ch00-tool-choice}
\begin{tabular}{p{0.36\linewidth}p{0.18\linewidth}p{0.36\linewidth}}
\toprule
確認したいこと & 主な道具 & 理由 \\
\midrule
代表的なユースケースで期待結果になる & テスト & 具体的な入出力例を読みやすく残せる \\
不正な形のデータを受け取らない & 型 & そもそも表現できない値を減らせる \\
業務上守るべきルールを明文化する & 仕様 & 人間が合意すべき性質を言葉にできる \\
すべての入力で式変形が意味を変えない & 証明 & 有限個の例ではなく一般的な主張を確認できる \\
旧形式から新形式へ移して戻せる & 証明 & round-trip のような全称的性質を直接表せる \\
外部 API が実際に応答する & テスト・監視 & モデル外の環境依存を観測する必要がある \\
\bottomrule
\end{tabular}
\end{table}

Lean での証明に適しているのは、「小さく切り出せる」「仕様として重要」「具体例では網羅しにくい」「変更時に壊れると困る」という条件を満たす性質です。

\section{圏論を合成と変換の言語として見る}
\label{sec:ch00-category-language}

\subsection{圏論の最初の読み替え}

圏論という名前は抽象的に見えます。
しかし、本書では最初から高度な一般論には入りません。
まずは、ソフトウェアで慣れた言葉へ読み替えます。

\begin{table}[htbp]
\centering
\caption{圏論とソフトウェアの語彙対応}
\label{tab:ch00-category-translation}
\begin{tabular}{p{0.20\linewidth}p{0.34\linewidth}p{0.36\linewidth}}
\toprule
圏論の語彙 & ソフトウェアでの読み替え & 例 \\
\midrule
対象 & 型、データ構造、状態空間 & \texttt{UserV1}, \texttt{Order}, \texttt{Account} \\
射 & 関数、変換、処理 & \texttt{normalize}, \texttt{migrate}, \texttt{decode} \\
合成 & パイプライン、処理連結 & \texttt{validate} の後に \texttt{persist} する \\
恒等射 & 何もしない処理 & 変換層を通しても値を変えない \\
同型 & 情報を失わない相互変換 & \texttt{migrate} と \texttt{rollback} が戻し合う \\
関手 & 構造を保つ一括変換 & \texttt{List.map migrate} \\
自然変換 & 一貫した adapter & \texttt{OldResponse A} から \texttt{NewResponse A} への変換 \\
積 & 両方を持つ構造 & ペア、レコード、DTO \\
余積 & どちらか一方の構造 & \texttt{Sum}, \texttt{Either}, \texttt{Result} \\
モナド & 文脈つき処理の合成 & 失敗、状態、検証を含む処理 \\
\bottomrule
\end{tabular}
\end{table}

圏論を使うということは、すべてを数学的な表現へ言い換えることではありません。
実務上の変換を、合成できるか、戻せるか、構造を保存するか、一貫しているか、という観点で整理することです。

\subsection{変更を矢印として見る}

本書では、処理や変更をしばしば矢印として描きます。

\[
  \texttt{UserV1} \xrightarrow{\texttt{migrate}} \texttt{UserV2}
\]

この図は、単に「旧形式から新形式へ変換する」という意味だけを表すものではありません。
次の問いを引き出すための図です。

\begin{itemize}
\item 逆向きの変換はありますか。
\item 変換して戻したら元に戻りますか。
\item 逆向きにも元に戻りますか。
\item すべての値で成り立ちますか、それとも well-formed な値だけで成り立ちますか。
\item リスト全体へ \texttt{map} しても件数や順序は保たれますか。
\item adapter を通す前後で業務ロジックの結果は変わりませんか。
\end{itemize}

図\ref{fig:ch00-change-map}は、旧形式と新形式、正逆の変換、変換後にも保存したい性質を別々の要素として配置します。

\FigurePlaceholder{fig-ch00-change-map.png}{0.90}{変更と保存性}{fig:ch00-change-map}

このように捉えると、圏論の語彙は実務レビューの語彙になります。
「これは同型なのか」「これは片方向互換なのか」「この \texttt{map} は構造を保存しているのか」「この adapter は自然なのか」という問いは、抽象数学のためだけのものではありません。
安全な変更を議論するための問いでもあります。

\section{Lean 4 を仕様を検査する小さな実験室として使う}
\label{sec:ch00-lean-lab}

\subsection{Lean で何をするのか}

Lean 4 は、プログラムを書ける言語であり、命題と証明も書ける言語です。
本書では、この二面性を利用します。

たとえば、\texttt{UserV1} と \texttt{UserV2} という二つの小さな構造体を Lean で定義します。
次に、\texttt{migrate : UserV1 \to UserV2} と \texttt{rollback : UserV2 \to UserV1} を書きます。
さらに、\texttt{rollback (migrate u) = u} を定理として書きます。
Lean がその証明を受け入れれば、少なくともその小さなモデルでは、移行して戻すと元に戻ることを確認できます。

Lean で行うのは、「本番コード全体をそのまま証明する」ことではありません。
まず、仕様の核を小さな型、関数、命題として表します。
そのうえで、重要な性質がすべての入力に対して成り立つことを確認します。

次のコードは、現時点で読めなくても問題ありません。
細部の構文は、\ChapterRef{chap:ch01_reading_code}以降で一つずつ説明します。
本書を読み終えるころには、この形の主張を自分で読めるようになります。
ここでは、\texttt{migrate} と \texttt{rollback} という二つの変換があり、\texttt{rollback (migrate u) = u}（移行して戻すと元に戻る）という性質を証明したい、という点だけを確認します。

\begin{listing}[htbp]
\begin{minted}{lean4}
structure UserV1 where
  id : Nat
  name : String

deriving instance Repr for UserV1

structure UserV2 where
  id : Nat
  profileName : String

deriving instance Repr for UserV2

def migrate (u : UserV1) : UserV2 :=
  { id := u.id, profileName := u.name }

def rollback (v : UserV2) : UserV1 :=
  { id := v.id, name := v.profileName }

theorem rollback_migrate (u : UserV1) :
    rollback (migrate u) = u := by
  cases u
  rfl

theorem migrate_rollback (v : UserV2) :
    migrate (rollback v) = v := by
  cases v
  rfl
\end{minted}
\caption{\texttt{migrate}・\texttt{rollback}・round-trip}
\label{lst:ch00_roundtrip_preview}
\end{listing}

この例の要点は、Lean の構文を覚えることではありません。
\texttt{migrate} と \texttt{rollback} があるだけでは不十分であり、二つの round-trip 性質を確認して初めて「情報を失っていない」と言えます。

\subsection{証明モデルと本番実装の境界}

Lean で証明した性質は、Lean で定義したモデルの範囲で強い保証を持ちます。
現実のシステムには、モデル化していない要素が残ります。

Lean で \texttt{rollback (migrate u) = u} を証明しても、実装言語の JSON ライブラリ、DB の NULL 制約、文字コード、タイムゾーン、外部 API の仕様変更まで自動で保証されるわけではありません。
証明した範囲と、テストやレビューで確認すべき範囲を分けて書く必要があります。

したがって、本書では次の方針を採ります。

\begin{itemize}
\item 実務の仕様から、証明したい小さな性質を切り出します。
\item その性質を Lean の型、関数、述語、定理として表します。
\item 証明が通る範囲を明記します。
\item モデル外の要素は、通常のテスト、レビュー、監視と組み合わせます。
\end{itemize}

この方針によって、Lean の保証範囲を過大評価せずに済みます。
Lean は厳密な道具ですが、モデル化した範囲を超えることは証明しません。

\section{本書で扱うサンプル一覧}
\label{sec:ch00-samples}

本書では、抽象概念を先に積み上げるのではなく、ソフトウェア上の問題から出発します。
次のサンプルは、章をまたいで繰り返し登場します。

図\ref{fig:ch00-case-studies}は、各ケーススタディを、意味、情報、不変条件、クエリ結果などの保存対象へ結びつけます。

\FigurePlaceholder{fig-ch00-case-studies.png}{0.92}{ケーススタディと保存対象}{fig:ch00-case-studies}

\begin{table}[htbp]
\centering
\caption{主なケーススタディ}
\label{tab:ch00-samples}
\begin{tabular}{p{0.28\linewidth}p{0.32\linewidth}p{0.30\linewidth}}
\toprule
サンプル & 保存したい性質 & 関連する主な章 \\
\midrule
関数パイプラインのリファクタリング & 入力に対する出力が変わらない & \ChapterRef{chap:ch09_equality_refactoring}、\ChapterRef{chap:ch11_category}、\ChapterRef{chap:ch18_equations_rewriting_normalization}、\ChapterRef{chap:ch34-refactoring-as-equations} \\
\texttt{UserV1} から \texttt{UserV2} へのスキーマ移行 & round-trip によって情報が戻る & \ChapterRef{chap:ch12-isomorphism}、\ChapterRef{chap:ch27-migration-classification}、\ChapterRef{chap:ch28_user_schema_lossless_migration} \\
JSON encoder / decoder & encode して decode すると元に戻る & \ChapterRef{chap:ch12-isomorphism}、\ChapterRef{chap:ch27-migration-classification} \\
口座残高操作 & 操作後も不変条件が保たれる & \ChapterRef{chap:ch10_invariants}、\ChapterRef{chap:spec-to-lean}、\ChapterRef{chap:account-spec-proof} \\
API adapter & adapter と業務ロジックの順序を入れ替えても結果が一致する & \ChapterRef{chap:natural-transformation}、\ChapterRef{chap:function-extensionality}、\ChapterRef{chap:ch31-api-adapter-naturality} \\
DB join と整合性制約 & 移行後も必要な join 結果が保存される & \ChapterRef{chap:limits-pullbacks}、\ChapterRef{chap:db-schema-query-preservation} \\
\bottomrule
\end{tabular}
\end{table}

\subsection{関数パイプラインのリファクタリング}

リファクタリングでは、内部構造を変えても外から見た意味を変えないことが求められます。
Lean では、この主張を等式として書きます。

\[
  \texttt{newPipeline x = oldPipeline x}
\]

この形は単純に見えますが、任意の \texttt{x} についての強い主張です。
この等式が成り立てば、特定のテストケースだけでなく、モデル化されたすべての入力で意味が保存されます。

\subsection{スキーマ移行と round-trip}

\texttt{UserV1} から \texttt{UserV2} への移行では、形が変わっても情報が失われないかを確認します。
情報を完全に保つ場合は、次の二つの式を証明する必要があります。

\[
  \texttt{rollback (migrate u) = u}
\]

\[
  \texttt{migrate (rollback v) = v}
\]

片方だけが成り立つ場合は、完全同型ではなく片方向互換である可能性があります。
条件付きでしか成り立たない場合は、well-formed 条件を仕様に入れる必要があります。

\subsection{JSON encoder / decoder}

encoder と decoder も、同型や round-trip の典型例です。
ただし、実務では、空白、フィールド順、デフォルト値、省略フィールド、未知フィールド、数値表現などを考慮する必要があります。
これらの要素によって、「文字列として同じ」と「データとして同じ」が一致しないことがあります。

JSON の例では、何を同じと見なすかを仕様として切り出します。
データとして戻ればよいのか、表示文字列まで一致させる必要があるのかを区別します。

\subsection{口座残高操作}

口座操作では、入金、出金、送金を扱います。
中心となるのは不変条件です。
たとえば、「残高は負にならない」という条件を操作後にも保つ必要があります。

この性質は、実務ではテストしやすいように見えます。
しかし、境界値や失敗ケースを含むすべての入力を手作業で列挙することは困難です。
Lean では、モデル化した範囲で「すべての入力」に対する主張として扱えます。

\subsection{API adapter}

API レスポンスの wrapper を変更するとき、旧形式から新形式へ変換する adapter を作ることがあります。
このとき、単に adapter が存在するだけでは不十分です。
業務ロジックを適用してから adapter を通す場合と、adapter を通してから業務ロジックを適用する場合で結果が一致するなら、adapter は業務ロジックと干渉していないと言えます。

この性質は自然性であり、実務的には「どの型でも一貫して動く adapter の条件」と読めます。

\subsection{DB join と整合性制約}

DB スキーマ変更では、移行前後でテーブルの形が変わります。
保存したいものは、必ずしもテーブル構造そのものではありません。
業務上必要なクエリ結果が保存されること、特に join の意味が変わらないことが求められます。

このとき、単に二つのレコードをペアにするだけでは足りません。
ユーザー ID が一致すること、注文が存在すること、参照整合性が満たされることなどの制約が必要です。
Pullback は、この整合条件を満たすペアを表す構造として読めます。

\section{本書の読み方}
\label{sec:ch00-reading-guide}

\subsection{最短実務ルート}

Lean 4 と圏論の両方が初めてで、まず実務に使える感覚を得たい場合は、次の順で読みます。

\begin{enumerate}
\item \ChapterRef{chap:ch01_reading_code}から\ChapterRef{chap:ch10_invariants}で、Lean 4 における型、関数、命題、不変条件を学びます。
\item \ChapterRef{chap:ch18_equations_rewriting_normalization}から\ChapterRef{chap:function-extensionality}で、等式、場合分け、帰納法、関数外延性という証明道具を身につけます。
\item \ChapterRef{chap:ch11_category}から\ChapterRef{chap:natural-transformation}で、圏、同型、積・余積、関手、自然変換をソフトウェアの例で理解します。
\item \ChapterRef{chap:spec-to-lean}から\ChapterRef{chap:ch31-api-adapter-naturality}で、仕様証明、マイグレーション、API adapter のケーススタディを確認します。
\item 付録のチェックリストを使い、自分のプロジェクトの変更に当てはめます。
\end{enumerate}

\subsection{Lean 証明力を先に強めるルート}

Lean で仕様を書けるようになりたい場合は、次の順が適しています。

\begin{enumerate}
\item \ChapterRef{chap:ch01_reading_code}から\ChapterRef{chap:ch10_invariants}で Lean の最小構文を学びます。
\item \ChapterRef{chap:ch18_equations_rewriting_normalization}から\ChapterRef{chap:function-extensionality}で、等式、場合分け、帰納法、関数外延性を学びます。
\item \ChapterRef{chap:spec-to-lean}から\ChapterRef{chap:ch26_price_spec}で、自然言語仕様を Lean の定理へ落とす流れを確認します。
\end{enumerate}

このルートでは、圏論の抽象語彙が必要になった時点で該当箇所へ戻ってもかまいません。

\subsection{圏論を丁寧に学ぶルート}

圏論の概念そのものを丁寧に身につけたい場合も、まず\ChapterRef{chap:ch01_reading_code}から\ChapterRef{chap:ch10_invariants}、続いて\ChapterRef{chap:ch18_equations_rewriting_normalization}から\ChapterRef{chap:function-extensionality}で Lean の証明道具を確認します。
そのうえで、\ChapterRef{chap:ch11_category}から\ChapterRef{chap:monad}をゆっくり読みます。
対象、射、合成、同型、関手、自然変換、モノイド、モナドを、すべてソフトウェアの例と結びつけて理解します。
\ChapterRef{chap:relations-preorders-galois}、\ChapterRef{chap:limits-pullbacks}、\ChapterRef{chap:ch37_adjunction}以降は、これらの概念をより抽象的な構造として整理します。

どのルートでも、自作の小さな定義による概念理解と、Mathlib の \texttt{CategoryTheory} API による標準化を分けて扱います。
この分離により、Lean と圏論を同時に学ぶ負荷を減らせます。

\section{本章のまとめ}

ソフトウェアの正しさは、実行結果、型、仕様、不変条件、変更前後の意味保存など、複数の層に分かれます。

テスト、型、仕様、証明は互いに補完する道具であり、どれか一つがほかを完全に置き換えるわけではありません。

圏論は、型やデータ構造を対象、関数や変換を射、パイプラインを合成として見るための言語です。

Lean 4 は、仕様の核を小さなモデルとして表し、すべての場合に成り立つ性質を機械的に確認するための実験室です。

本書では、リファクタリング、スキーマ移行、encoder / decoder、口座操作、API adapter、DB join を通じて、「保存されるべき性質」を具体的に証明します。
